% БҮЛЭГ 5 — ЗАГВАР ХӨГЖҮҮЛЭЛТ БА PIPELINE

\chapter{Загвар хөгжүүлэлт ба Pipeline}
\label{Chapter5}

Энэхүү бүлэгт бүлэг~\ref{Chapter3}-т тайлбарласан архитектурын дагуу MN-BERT
суурьтай хоёр шатлалт системийн загварын бүтэц, сургалтын pipeline, гиперпараметрийн
оновчлол, ангиллын тэнцвэржүүлэлт болон туршилтын давтагдах байдлын (reproducibility)
нөхцөлийг нарийвчлан тайлбарлана. Уг бүлгийн зорилго нь загварын хөгжүүлэлтийн
бүхий л үе шатыг баримтжуулж, дараагийн бүлэгт (бүлэг~\ref{Chapter6}) танилцуулагдах
туршилтын үр дүнг үнэмшилтэйгээр тайлбарлахад оршино.

%-------------------------------------------------------------------------------
\section{MN-BERT загварын архитектурын дэлгэрэнгүй тайлбар}

MN-BERT нь Монгол хэлний corpus дээр pre-train хийгдсэн BERT--base архитектур
бөгөөд \texttt{bert-base-mongolian-cased} хувилбарыг суурь загвар болгон
ашигласан. Уг загвар нь 12 Transformer энкодер давхарга (encoder layer), 12
attention head, далд давхаргын хэмжээ 768, нийт ойролцоогоор 110 сая параметртэй.
Оролтын хамгийн их урт нь 512 token бөгөөд сентимент ангилалд 128 token-ы урт
ашиглахаар тохируулсан --- учир нь бүлэг~\ref{Chapter2}-т танилцуулсан өгөгдлийн
ихэнх зүйл энэ хязгаарт багтдаг.

Ангиллын даалгаварт зориулан MN-BERT-ийн дээр дараах давхаргуудыг нэмсэн:

\begin{enumerate}
  \item \textbf{Tokenization давхарга:} WordPiece-т суурилсан subword tokenizer
    нь Монгол хэлний agglutinative морфологийг үлдэгдэл бус (lossless) байдлаар
    кодолно. Оролтын текст эхэнд \texttt{[CLS]}, төгсгөлд \texttt{[SEP]} тэмдэгт
    нэмэгдэнэ.
  \item \textbf{Transformer encoder stack:} 12 давхаргын гаралт нь контекст мэдрэмжтэй
    vector илэрхийлэл үүсгэнэ. Энэ хэсэг нь pre-train хийгдсэн жингүүдийг үндэс
    болгон fine-tuning хийгдэнэ.
  \item \textbf{[CLS] pooling давхарга:} Төгсгөлийн encoder давхаргын
    \texttt{[CLS]} байрлалын далд вектор нь өгүүлбэрийн нэгдсэн дүрслэл болж
    ангиллын толгойд (classification head) дамжина.
  \item \textbf{Ангиллын толгой (classification head):} Ганц бүрэн холбосон
    (fully connected) давхарга нь pooled vector-ыг ангиллын логитууд (logits)
    руу шугаман хувиргана. Stage~1-ийн хувьд гаралтын хэмжээ 3 (Positive,
    Neutral, Negative), Stage~2-ийн хувьд 2 (Toxic, Constructive).
  \item \textbf{Softmax давхарга:} Логитуудыг магадлалын тархалтад хувиргана.
    Инференс (inference) үедээ \texttt{argmax} ашиглан эцсийн ангиллыг тодорхойлно.
\end{enumerate}

Stage~1 болон Stage~2 нь тус тусдаа fine-tune хийгдсэн бие даасан загварууд
бөгөөд мэдээлэл солилцох бус, каскад (cascade) хэлбэрээр холбогдсон. Энэ нь
ангилал тус бүрийг өөрийн ангиллын бүрэлдэхүүнд оновчтой дасан зохицохыг
боломжтой болгож байна.

%-------------------------------------------------------------------------------
\section{Stage 1 ба Stage 2 тус тусын training pipeline}

Stage~1 болон Stage~2 хоёулаа PyTorch 2.x framework, HuggingFace Transformers
4.x ашиглан сургагдсан. Сургалтын үйл явцыг нэгдмэл pipeline-д зохион байгуулсан
бөгөөд түүнийг хүснэгт~\ref{tab:pipeline-flow}-д багцалсан.

\begin{table}[H]
\centering
\caption{Сургалтын pipeline-ын үе шатуудын урсгал}
\label{tab:pipeline-flow}
\begin{tabular}{@{}clp{8cm}@{}}
\toprule
\textbf{№} & \textbf{Үе шат} & \textbf{Тодорхойлолт} \\
\midrule
1 & Өгөгдөл ачаалалт       & CSV файлаас pandas ашиглан унших, train/val/test хуваалтаар ангилах \\
2 & Tokenization           & \texttt{AutoTokenizer.from\_pretrained(''bert-base-mongolian-cased'')}, max\_length=128 \\
3 & PyTorch Dataset        & \texttt{input\_ids}, \texttt{attention\_mask}, \texttt{labels} tensor-уудыг боловсруулах \\
4 & DataLoader             & Batch болгон ачаалах, \texttt{shuffle=True} train-д, \texttt{shuffle=False} val/test-д \\
5 & Forward pass           & \texttt{outputs = model(input\_ids, attention\_mask, labels)}, cross-entropy loss тооцох \\
6 & Backward pass          & \texttt{loss.backward()} ба градиентын clipping (\texttt{max\_norm=1.0}) \\
7 & Optimizer step         & AdamW optimizer, weight decay зохицуулалттай \\
8 & LR scheduler           & Linear warmup (10\% алхамд), дараа нь шугаман буурах \\
9 & Валидаци               & Epoch тутам validation set-ээр F1-macro үнэлэх, \texttt{best} жинг хадгалах \\
10 & Чекпойнт              & \texttt{best\_model.pt} файлд жинг хадгалах (early stopping-тай) \\
\bottomrule
\end{tabular}
\end{table}

Stage~1-ийн хувьд оролтын өгөгдөл нь Positive / Neutral / Negative гэсэн гурван
ангилалтай, гаралт нь 3 логиттой cross-entropy loss. Stage~2-ийн хувьд зөвхөн
Negative гэж тэмдэглэгдсэн жишээнүүд ашиглагдах ба оролт нь хоёртын (binary)
ангилалтай (Toxic ба Constructive). Хоёр шат нь өөр өөр өгөгдлийн багц дээр
сургагдах тул тус тусын ``best'' жингээ хадгалж, инференс үед цувуулан дуудна.

%-------------------------------------------------------------------------------
\section{Оновчтой шийдэл: Нэг шатлалт (Flat) загварт шилжсэн нь}

Хоёр шатлалт каскад архитектурыг турших явцад гарсан үр дүн болон алдааны
дүн шинжилгээ (бүлэг \ref{Chapter6}-д дэлгэрэнгүй үзүүлсэн) нь нэгэн чухал
инженерийн сорилтыг илрүүлсэн. Энэ нь \textbf{Алдааны дамжуулалт (Error
Propagation)} буюу Stage~1 дээр гарсан өчүүхэн алдаа Stage~2-ын үр дүнг
бүрэн боломжгүй болгох bottleneck үүсгэсэн явдал юм.

Энэхүү олдвор дээр үндэслэн бид судалгааны чиглэлээ өөрчилж (pivot), дөрвөн
ангиллыг нэгэн зэрэг таамаглах \textbf{Нэг шатлалт (Flat) MN-BERT} загварыг
оновчтой шийдэл болгон сонгосон. Flat загварын давуу тал нь:

\begin{enumerate}
    \item \textbf{Алдааны нийлбэрээс сэргийлэх:} Бүх ангиллын хил хязгаарыг
    нэгэн зэрэг сурснаар Stage~1-ийн хатуу шүүлтүүрээс хамааралгүй болно.
    \item \textbf{Focal Loss ашиглалт:} Классуудын тэнцвэргүй байдлыг шийдэхэд
    Cross-Entropy-ийн оронд Focal Loss ашигласнаар хүнд жишээнүүд дээрх (Toxic)
    загварын төвлөрлийг нэмэгдүүлсэн.
    \item \textbf{Инференс хурд:} Хоёр загвар цувуулан ачаалах шаардлагагүй тул
    хариу өгөх хугацааг 2 дахин бууруулсан.
\end{enumerate}

Ийнхүү бидний эцсийн систем нь Focal Loss болон Cosine Warmup тохируулга бүхий
Flat MN-BERT загвар дээр суурилсан бөгөөд энэ нь туршилтын хамгийн өндөр
гүйцэтгэлийг үзүүлсэн юм.

%-------------------------------------------------------------------------------
\section{Гиперпараметрийн тохиргоо ба оновчлол}

Гиперпараметрийн сонголтыг бүлэг~\ref{Chapter3}-т танилцуулсан тохиргооны
үндсэн дээр хийж, Stage~1 ба Stage~2-д тус тусад нь validation F1-macro
метрикээр сонгож оновчлов. Эцсийн сонголтыг хүснэгт~\ref{tab:hyperparams-mnbert}-д
жагсаав.

\begin{table}[H]
\centering
\caption{MN-BERT fine-tuning-ын гиперпараметрийн эцсийн тохиргоо}
\label{tab:hyperparams-mnbert}
\begin{tabular}{@{}lcc@{}}
\toprule
\textbf{Гиперпараметр} & \textbf{Stage 1} & \textbf{Stage 2} \\
\midrule
Learning rate              & $3 \times 10^{-5}$ & $2 \times 10^{-5}$ \\
Batch size                 & 16                 & 16                 \\
Epochs (Best)              & 6                  & 5                  \\
Warmup ratio               & 0.1                & 0.1                \\
Weight decay               & 0.01               & 0.01               \\
Dropout                    & 0.1                & 0.1                \\
Max sequence length        & 128                & 128                \\
Optimizer                  & AdamW              & AdamW              \\
LR scheduler               & Cosine with Warmup & Cosine with Warmup \\
Early stopping patience    & 3                  & 3                  \\
\bottomrule
\end{tabular}
\end{table}

Learning rate-ийн хувьд $\{2\!\times\!10^{-5},\; 3\!\times\!10^{-5},\; 5\!\times\!10^{-5}\}$
гэсэн утгуудыг grid search аргаар туршин үзэж, хамгийн сайн validation F1-ыг
өгсөн утгыг сонгов. Batch size-ыг $\{16, 32\}$ хүрээнд туршихад 12GB VRAM-ийн
хязгаарлалтаас шалтгаалан batch size 16 нь хамгийн тогтвортой гүйцэтгэл үзүүлсэн.

%-------------------------------------------------------------------------------
\section{Сургалтын дэд бүтэц ба орчин}

Загваруудыг сургах болон турших ажлыг локал тооцооллын орчинд гүйцэтгэсэн.
Техник хангамжийн хувьд \textbf{AMD Radeon RX 9070 GRE 12GB} график картыг
ашигласан бөгөөд \texttt{torch-directml} санг ашиглан DirectX-ээр дамжуулан
GPU тооцооллыг идэвхжүүлсэн.

\begin{itemize}
    \item \textbf{GPU:} AMD Radeon RX 9070 GRE (12GB VRAM)
    \item \textbf{CPU:} Intel Core i7-12700K
    \item \textbf{RAM:} 32GB DDR5
    \item \textbf{Framework:} PyTorch 2.4.1 (with DirectML)
    \item \textbf{OS:} Windows 11
\end{itemize}

%-------------------------------------------------------------------------------
\section{SMOTE хэрэглэлт ба ангиллын тэнцвэржүүлэлт}

Бүлэг~\ref{Chapter2}-т дурдсанчлан, өгөгдлийн ангиллын тархалт нь тэнцвэргүй
байдлыг илэрхийлдэг: Neutral болон Positive ангилал давамгайлж, Constructive
ангилал хамгийн цөөхөн дээжтэй. Уг тэнцвэргүй байдал нь цөөнх ангиллын F1
гүйцэтгэлийг доройтуулж болзошгүй тул хоёр түвшинд тэнцвэржүүлэлт хийв.

\textbf{Baseline загваруудын хувьд (LR, NB, SVM):} \texttt{imbalanced-learn}
сангийн \texttt{SMOTE} (Synthetic Minority Over-sampling Technique) алгоритмыг
TF-IDF vector-ын дээр хэрэглэв. SMOTE нь цөөнх ангиллын жишээг $k$-хамгийн
ойрын хөршийнх нь интерполяциар синтетикээр үржүүлнэ. Хэрэгжүүлэлтэд
$k\!=\!5$ утгыг ашиглав.

\textbf{MN-BERT загварын хувьд:} BERT-д шууд SMOTE хэрэглэх нь эмбэддингийн
орон зайд семантикийг алдагдуулах эрсдэлтэй тул ангиллын жин тооцсон
cross-entropy loss-ыг ашиглав:
\[
  w_c = \frac{N}{K \cdot n_c}
\]
энд $N$ --- нийт жишээний тоо, $K$ --- ангиллын тоо, $n_c$ --- $c$-р
ангиллын жишээний тоо. Уг жинг PyTorch-ийн \texttt{CrossEntropyLoss(weight=...)}
параметрт дамжуулснаар цөөнх ангиллын алдаанд илүү жин өгнө.

Энэ хоёр арга нь бүлэг~\ref{Chapter6}-д үзүүлэх F1-macro-ийн сайжруулалтыг
өгөх ба guideline шаардлагын НШ-5 (``ангиллын тэнцвэржүүлэлт нь F1 сайжруулалтыг
үзүүлэх ёстой'') зорилгыг хангана.

%-------------------------------------------------------------------------------
\section{Gradio demo pipeline integration}

Эцсийн системийн хэрэглэгчийн интерфейсийг Gradio framework дээр бүтээсэн
бөгөөд pipeline-ийн интеграцийг дараах логикоор зохион байгуулсан:

\begin{enumerate}
  \item \textbf{Оролт:} Хэрэглэгч Монгол хэлний текстийг оруулна.
  \item \textbf{Урьдчилсан боловсруулалт:} Бүлэг~\ref{Chapter2}-т тодорхойлсон
    pipeline (цэвэрлэгээ, кирилл шүүлт, tokenization)-аар дамжина.
  \item \textbf{Stage~1 инференс:} \texttt{best\_stage1.pt} ачаалагдсан MN-BERT
    загвар Positive / Neutral / Negative ангиллыг магадлалтай буцаана.
  \item \textbf{Нөхцөлт салаалалт:} Хэрэв таамагласан ангилал Negative биш бол
    үр дүнг шууд хэрэглэгчид буцаана. Negative бол Stage~2 руу дамжина.
  \item \textbf{Stage~2 инференс:} \texttt{best\_stage2.pt} ачаалагдсан загвар
    Toxic / Constructive ангиллыг тодорхойлно.
  \item \textbf{Гаралт:} Эцсийн ангилал, магадлалын утга, attention-д суурилсан
    гол үгсийн тодруулга (highlight) болон тайлбар хэрэглэгчид харагдана.
\end{enumerate}

Загваруудыг эхний удаад санах ойд ачаалж (model loading), дараагийн хүсэлтүүдэд
re-load-гүй хурдан хариу өгөх боломжтой болгосон. GPU орчинд (AMD DirectML)
инференс дундаж хариултын хугацаа 20-35 миллисекунд, харин CPU орчинд
200-300 миллисекунд байна.

\begin{figure}[H]
    \centering
    \begin{tikzpicture}[
        node distance=1.2cm and 2cm,
        startstop/.style={rectangle, rounded corners, minimum width=4cm, minimum height=1.1cm, text centered, draw=black, fill=red!10, drop shadow, font=\small, align=center},
        process/.style={rectangle, minimum width=4cm, minimum height=1.1cm, text centered, draw=black, fill=orange!10, drop shadow, font=\small, align=center},
        decision/.style={diamond, aspect=2, minimum width=3.5cm, text centered, draw=black, fill=green!10, drop shadow, font=\small, inner sep=0pt, align=center},
        arrow/.style={thick, ->, >=stealth}
    ]
    \node (in) [startstop] {Оролт: Текст};
    \node (pre) [process, below=of in] {Цэвэрлэгээ ба \\ Tokenization};
    \node (s1) [process, below=of pre] {MN-BERT Stage 1 \\ (Sentiment)};
    \node (dec) [decision, below=1.5cm of s1] {Сөрөг? (Neg)};
    
    \node (s2) [process, right=2cm of dec] {MN-BERT Stage 2 \\ (Toxic/Con)};
    \node (out) [startstop, below=1.5cm of dec] {Эцсийн үр дүн};

    \draw [arrow] (in) -- (pre);
    \draw [arrow] (pre) -- (s1);
    \draw [arrow] (s1) -- (dec);
    \draw [arrow] (dec) -- node[anchor=south] {Тийм} (s2);
    \draw [arrow] (dec) -- node[anchor=west] {Үгүй} (out);
    \draw [arrow] (s2) |- (out);
    \end{tikzpicture}
    \caption{Системийн ажиллагааны логик ба инференс pipeline}
    \label{fig:inference-pipeline}
\end{figure}

%-------------------------------------------------------------------------------
\section{Reproducibility ба туршилтын бүртгэл}

Туршилтын давтагдах байдлыг (reproducibility) хангахын тулд дараах арга хэмжээг
авсан:

\begin{itemize}
  \item \textbf{Random seed тогтворжуулалт:} \texttt{torch.manual\_seed(42)},
    \texttt{numpy.random.seed(42)}, \texttt{random.seed(42)} бүх туршилтанд
    адилхан тохируулсан. CUDA deterministic горимыг
    (\texttt{torch.backends.cudnn.deterministic = True}) идэвхжүүлсэн.
  \item \textbf{Хувилбарын хяналт:} Бүх кодыг Git/GitHub repository-д хадгалж,
    commit hash бүрээр туршилтыг хамтад нь тэмдэглэв. \texttt{requirements.txt}
    файлд бүх сангийн яг тохирох хувилбарыг бичсэн.
  \item \textbf{Туршилтын бүртгэл:} Epoch тутмын loss, accuracy, F1 утгуудыг
    JSON лог файлд хадгалав. Сургалтын муруй (learning curve) болон
    confusion matrix-ыг \texttt{matplotlib}/\texttt{seaborn}-оор хадгалж,
    бүлэг~\ref{Chapter6}-д танилцуулна.
  \item \textbf{Загварын чекпойнт:} Хамгийн сайн validation F1-тэй epoch-ийн
    жинг \texttt{best\_stage1.pt}, \texttt{best\_stage2.pt} файлуудад хадгалав.
  \item \textbf{Орчны тохиргоо:} Google Colab дээр ажиллуулсан notebook-уудыг
    PDF болгон экспортлов. NVIDIA T4 болон A100 GPU-гийн хоёр төрөл дээр
    давтан сургаж, үр дүнгийн тогтвортой байдлыг баталгаажуулав.
\end{itemize}

Эдгээр арга хэмжээ нь guideline-ийн Ш-6 (``reproducible experiment setup
documented'') шаардлагыг бүрэн хангаж байна.

%-------------------------------------------------------------------------------
\section{Бүлгийн дүгнэлт}

Энэ бүлэгт MN-BERT суурьтай Stage~1 ба Stage~2 загваруудын архитектур, сургалтын
pipeline-ийн бүх алхмууд, гиперпараметрийн сонголт, ангиллын тэнцвэржүүлэлтийн
хоёр арга (SMOTE baseline-д, жингэсэн cross-entropy BERT-д), Gradio-д
суурилсан инференсийн урсгал, болон reproducibility-ийн нөхцөлийг тодорхой
баримтжуулсан. Сургалтын бүх бүрэлдэхүүн хэсэг нь нийтлэг Python/PyTorch
экосистем дээр суурилсан тул бусад судлаачид давтан гүйцэтгэх боломжтой. Ийнхүү
загварын хөгжүүлэлтийн үе шат бэлэн болсон бөгөөд дараагийн бүлэгт уг
загваруудын туршилтын үр дүнг нарийвчлан авч үзнэ.
